% =========================================================
% Common Errors and Fallacies
% =========================================================

\subsection*{Common Errors and Fallacies}

<<<<<<< HEAD
\begin{toolkitbox}{Predicate Logic Error Toolkit --- Quick Reference}
\small
\begin{tabular}{p{0.28\linewidth} p{0.34\linewidth} p{0.28\linewidth}}
\toprule
\textbf{Check} & \textbf{Question} & \textbf{Typical repair} \\
\midrule
Quantifier choice & Is the claim about every object or about at least one object? & Use $\forall$ for all, $\exists$ for some. \\
Conditional translation & Is membership in a class a hypothesis? & Use implication for universal class claims. \\
Witness discipline & Has the chosen object been introduced legally? & Keep witnesses new and local. \\
Scope and substitution & Which variables are bound, free, or captured? & Rename bound variables before substituting. \\
\bottomrule
\end{tabular}
\end{toolkitbox}

\begin{remark*}[Orientation]
The mistakes in this section are not new theorem schemas. They are reference
checks for reading, translating, and proving with predicate logic. Before
correcting a formula, identify which variables are bound, which variables are
free, and which connective or quantifier has the largest scope.
\end{remark*}

\begin{remark}[Universal and existential quantifiers]\label{rem:predicate-logic-universal-existential-confusion}
\textbf{Incorrect form:} $\forall x\,P(x)$ when the intended claim is that
some object has property $P$.

\textbf{Corrected form:} $\exists x\,P(x)$.

\textbf{Explanation:} A universal statement commits to every object in the
domain; an existential statement commits only to at least one object.

\textbf{Example:} ``Some integer is prime'' is
$\exists n\,(Integer(n)\wedge Prime(n))$, not
$\forall n\,(Integer(n)\rightarrow Prime(n))$.
\end{remark}

\begin{remark}[Translating all with conjunction]\label{rem:predicate-logic-all-conjunction}
\textbf{Incorrect form:} $\forall x\,(A(x)\wedge B(x))$ for ``all $A$ are
$B$.''

\textbf{Corrected form:} $\forall x\,(A(x)\rightarrow B(x))$.

\textbf{Explanation:} A universal class statement says that being an $A$ is a
sufficient condition for being a $B$; it does not say every object in the
domain is an $A$.

\textbf{Example:} ``All squares are rectangles'' is
$\forall x\,(Square(x)\rightarrow Rectangle(x))$.
\end{remark}

\begin{remark}[Translating some with implication]\label{rem:predicate-logic-some-implication}
\textbf{Incorrect form:} $\exists x\,(A(x)\rightarrow B(x))$ for ``some $A$
are $B$.''

\textbf{Corrected form:} $\exists x\,(A(x)\wedge B(x))$.

\textbf{Explanation:} An existential class claim needs one object satisfying
both conditions. The implication may be true for an object that is not an $A$.

\textbf{Example:} ``Some integers are even'' is
$\exists n\,(Integer(n)\wedge Even(n))$.
\end{remark}

\begin{remark}[Quantifier scope errors]\label{rem:predicate-logic-scope-errors}
\textbf{Incorrect form:} $\forall x\,P(x)\rightarrow Q(x)$ when the intended
claim is that every $P$ is $Q$.

\textbf{Corrected form:} $\forall x\,(P(x)\rightarrow Q(x))$.

\textbf{Explanation:} Parentheses determine how far the quantifier reaches.
Without the intended scope, a variable may occur outside the formula governed
by its quantifier.

\textbf{Example:} ``Every differentiable function is continuous'' should bind
the whole implication:
$\forall f\,(Differentiable(f)\rightarrow Continuous(f))$.
\end{remark}

\begin{remark}[Invalid quantifier reversal]\label{rem:predicate-logic-invalid-quantifier-reversal}
\textbf{Incorrect form:}
\[
\forall x\exists y\,R(x,y)
\equiv
\exists y\forall x\,R(x,y).
\]

\textbf{Corrected form:} In general,
\[
\forall x\exists y\,R(x,y)
\not\equiv
\exists y\forall x\,R(x,y).
\]

\textbf{Explanation:} The first statement allows the witness $y$ to depend on
$x$. The second requires one single $y$ that works for every $x$.

\textbf{Example:} Every real number has an additive inverse, but there is no
single real number that is the additive inverse of every real number.
\end{remark}

\begin{remark}[Free variables in intended sentences]\label{rem:predicate-logic-free-variable-sentence}
\textbf{Incorrect form:} $P(x)\rightarrow Q(x)$ as a complete sentence.

\textbf{Corrected form:} $\forall x\,(P(x)\rightarrow Q(x))$, or another
explicit quantification matching the intended claim.

\textbf{Explanation:} A sentence has no free variables. A formula with a free
variable still depends on an assignment or context.

\textbf{Example:} ``Every prime greater than $2$ is odd'' is
$\forall n\,((Prime(n)\wedge n>2)\rightarrow Odd(n))$.
\end{remark}

\begin{remark}[Variable capture during substitution]\label{rem:predicate-logic-variable-capture}
\textbf{Incorrect form:} Substituting $y$ for $x$ in
$\exists y\,R(x,y)$ and writing $\exists y\,R(y,y)$.

\textbf{Corrected form:} First rename the bound variable:
$\exists z\,R(x,z)$, then substitute to get $\exists z\,R(y,z)$.

\textbf{Explanation:} A substitution is valid only when the substituted term is
free for the variable being replaced. Capture changes a free occurrence into a
bound occurrence.

\textbf{Example:} In $\forall y\,(x<y)$, substituting $y$ for $x$ without
renaming would falsely turn the free $y$ into a bound variable.
\end{remark}

\begin{remark}[Reusing a witness incorrectly]\label{rem:predicate-logic-reused-witness}
\textbf{Incorrect form:} From $\forall x\exists y\,R(x,y)$, choose the same
constant $c$ for every $x$.

\textbf{Corrected form:} For an arbitrary $x$, choose a witness that may
depend on that $x$; do not reuse it outside its allowed scope.

\textbf{Explanation:} Existential witnesses introduced under a universal
choice usually depend on that choice.

\textbf{Example:} From ``every person has a mother,'' one may choose a mother
for each person, but not a single person who is everyone's mother.
\end{remark}

\begin{remark}[Illicit universal generalization]\label{rem:predicate-logic-illicit-universal-generalization}
\textbf{Incorrect form:} Proving $P(a)$ for a special object $a$ and
concluding $\forall x\,P(x)$.

\textbf{Corrected form:} Conclude $\forall x\,P(x)$ only after proving
$P(x)$ for an arbitrary object independent of special assumptions.

\textbf{Explanation:} Universal generalization requires that the object being
generalized over not depend on an undischarged special condition.

\textbf{Example:} Showing that $0+0=0$ does not prove
$\forall n\,(n+n=n)$.
\end{remark}

\begin{remark}[Illicit existential instantiation]\label{rem:predicate-logic-illicit-existential-instantiation}
\textbf{Incorrect form:} From $\exists x\,P(x)$, introduce a witness $a$ that
is already used elsewhere with extra assumptions.

\textbf{Corrected form:} Introduce a new temporary constant $c$ with $P(c)$,
and keep conclusions independent of the particular name $c$.

\textbf{Explanation:} Existential instantiation names an unknown witness. The
name must not carry prior information.

\textbf{Example:} From ``some integer is even,'' one may introduce a fresh
witness $k$ with $Even(k)$, but not assume that the witness is the previously
chosen integer $n$.
\end{remark}

\begin{remark}[Affirming the consequent]\label{rem:predicate-logic-affirming-consequent}
\textbf{Incorrect form:} From $P\rightarrow Q$ and $Q$, infer $P$.

\textbf{Corrected form:} From $P\rightarrow Q$ and $P$, infer $Q$.

\textbf{Explanation:} A consequence may have many possible causes; observing
the consequence does not identify the antecedent.

\textbf{Example:} ``If a number is divisible by $4$, then it is even'' and
``$6$ is even'' do not imply that $6$ is divisible by $4$.
\end{remark}

\begin{remark}[Denying the antecedent]\label{rem:predicate-logic-denying-antecedent}
\textbf{Incorrect form:} From $P\rightarrow Q$ and $\neg P$, infer $\neg Q$.

\textbf{Corrected form:} From $P\rightarrow Q$ and $\neg Q$, infer $\neg P$.

\textbf{Explanation:} The antecedent may be only one sufficient condition for
the consequent.

\textbf{Example:} ``If a function is differentiable, then it is continuous''
and ``this function is not differentiable'' do not imply that the function is
not continuous.
\end{remark}

\begin{example}[A compact diagnostic pass]\label{ex:predicate-logic-fallacy-diagnostic}
Suppose a proof starts with
\[
\forall x\exists y\,R(x,y)
\]
and then fixes one object $c$ such that $\forall x\,R(x,c)$. The diagnostic is
``invalid quantifier reversal'' together with ``reusing a witness
incorrectly.'' The original statement allows $y$ to depend on $x$; it does not
provide one uniform witness.
\end{example}
=======
\begin{remark*}[Translation and proof errors]
\renewcommand{\arraystretch}{1.35}
\begin{tabular}{p{0.25\linewidth} p{0.25\linewidth} p{0.25\linewidth} p{0.17\linewidth}}
\toprule
\textbf{Error} & \textbf{Incorrect form} & \textbf{Corrected form} & \textbf{Explanation} \\
\midrule
Confusing \(\forall\) and \(\exists\) &
\(\exists x\,P(x)\) for ``every \(x\) has \(P\)'' &
\(\forall x\,P(x)\) &
Existence does not give universality. \\
Universal conjunction &
\(\forall x\,(A(x)\land B(x))\) &
\(\forall x\,(A(x)\to B(x))\) &
The antecedent restricts the domain. \\
Existential implication &
\(\exists x\,(A(x)\to B(x))\) &
\(\exists x\,(A(x)\land B(x))\) &
The witness must be both \(A\) and \(B\). \\
Scope error &
\(\exists y\forall x\,R(x,y)\) &
\(\forall x\exists y\,R(x,y)\), when each \(x\) may have its own \(y\) &
Quantifier order controls dependence. \\
Variable capture &
\(\varphi[t/x]\) without checking bound variables &
Rename bound variables before substitution &
Substitution must not create unintended binding. \\
Free variable left in a sentence &
\(\forall x\,R(x,y)\) as a closed claim &
\(\forall x\forall y\,R(x,y)\), or bind \(y\) as intended &
A sentence has no free variables. \\
Affirming the consequent &
\(P\to Q,\;Q\;\therefore P\) &
Invalid; use a counterexample valuation or structure &
The consequent may hold for another reason. \\
Denying the antecedent &
\(P\to Q,\;\neg P\;\therefore\neg Q\) &
Invalid; use a counterexample valuation or structure &
Failure of \(P\) does not force failure of \(Q\). \\
Invalid quantifier reversal &
\(\forall x\exists y\,R(x,y)\Rightarrow \exists y\forall x\,R(x,y)\) &
Invalid without extra hypotheses &
Dependent witnesses need not be uniform. \\
\bottomrule
\end{tabular}
\end{remark*}

\begin{remark*}[Failure-mode checklist]
Before accepting a translation or inference, check the domain, each symbol
reading, every bound-variable scope, and whether the final formula is intended
to be a sentence. Then test suspicious implications by trying to build one
case where the premise is true and the conclusion false.
\end{remark*}
>>>>>>> 591349ed2a1982198e862953b73959c91a638d59
